\thispagestyle{plain}
\ctsection{De hac editione}

Hæc editio totum Phraseologiæ Wagnerianæ corpus --- articulos 11\,239
--- ex editione Belgica anni MDCCCLXXVIII repetit. Suus cuique articulo
servatur apparatus: interpretationes Gallicæ litteris
inclinatis, synonyma (\dictlabel{syn.}), contraria (signo \mbox{)(}\,
notata), adverbia (\dictlabel{adv.}), epitheta (\dictlabel{epith.}),
phrases evolvendæ (\dictlabel{phras.}), usus auctorum
(\dictlabel{usus}) cum acceptionibus numeratis, notæque interdum
adiectæ (\dictlabel{vulg.}, \dictlabel{prov.}, \dictlabel{rad.},
\dictlabel{differ.}). Asteriscus voces ab auctore pro non Ciceronianis
habitas distinguit; parentheses locutiones vulgares includunt.

Lectiones mendosæ, etiam in præfatione Latina atque indicibus,
emendatæ sunt. Lectiones priores una cum emendationum rationibus et
testimoniis in commentario separato referuntur; emendationes ex
conjectura factæ a lectionibus testimoniis confirmatis distinguuntur.
In præfatione \textit{tum quod} et \textit{suppeditarent} ex conjectura
recepta sunt. In articulo \textit{accio}, pro \textit{exultro} legitur
\textit{ultro}; in indice Gallico-Latino, omisso \textit{Affermo},
\textit{loco} servatur.

Quantitates vocalium, ubi satis certa testimonia suppetunt, correctæ
sunt. Ubi de quantitate vocalis auctoritates dissentiunt vel
testimonia non sufficiunt, signum dubium omittitur: vocalis sine
signo neque longa neque brevis eo ipso declaratur. In lemmatibus
\textit{gratuitus} et \textit{istiusmodi} quantitas usitatior exhibetur;
brevior tamen vocalis apud poetas admittitur (\textit{gratuitum},
Statius, \textit{Silvæ} 1.6.16; genitivus in \textit{-ius}, Allen et
Greenough, \S~603). Scriptura \textit{quatuor}, quam Gaffiot quoque
admittit, servatur.

Tres indices historici editionis anni MDCCCLXXVIII volumen claudunt,
ut in originali: index vocum barbararum cum Latinis earum
æquivalentibus, index vocum quæ in foro militari, civili sacroque
obtinent, et index Gallico-Latinus.

\medskip
\noindent\textit{De usu vocabulorum.}
\textit{Etesiæ} plerumque masculino genere dicuntur; femininum tamen
apud Hyginum et Isidorum invenitur (Gaffiot; Georges).
\textit{Planeta} apud auctores classicos masculinum est; apud scriptores
medii ævi etiam femininum reperitur (\textit{Dictionary of Medieval
Latin from British Sources}). Utriusque vocabuli genus femininum,
quod hæc editio tradit, servatum est.
\textit{Obsurdeo}, in indice Gallico-Latino positum, etiam in
vocabulario quod Joris van Gheel anno MDCLII conscripsit legitur;
pro \textit{obsurdesco} usurpatum retinetur.
Genitivus pluralis \textit{sudium}, in indice rerum militarium
servatus, apud Sidonium quoque legitur (\textit{Epistulæ} 7.1.2).

\medskip
\noindent\textit{De præfatione Gallica.}
Locutio \textit{facere sumptus}, quam præfatio Gallica reprehendit,
non est per se a Latinitate aliena: Cicero ipse \textit{sumptum faciat}
(\textit{De imperio Cn. Pompei} 39) et \textit{sumptum in rem militarem
facere} (\textit{Ad familiares} 12.30.4) scripsit. Auctoris judicium
historicum servatur, sed hac nota definitur.
\vspace{1.5pc}
